\textbf{Background.} \textcite{huang2025} built machine-learning models
for mortality in intensive care unit (ICU) patients with diabetic
neuropathy (DN) from MIMIC-IV. After LASSO selection of 12 of 55
candidate predictors, random forest (RF) outperformed XGBoost, support
vector machine and logistic regression (LR) (validation AUC 0.780).
However, inclusion used information from after the first ICU day, the
outcome was not consistently defined, and models were chosen on one
263-patient validation set with selection and tuning outside
cross-validation.

\textbf{Methods.} Keeping the benchmark's framework, we used MIMIC-IV
v3.1 to build a landmark model predicting at 24 h after ICU admission
from information then available. We retained patients without a
white-cell differential, separated chronic from acute comorbidity codes
and defined death within 365 days. Models were evaluated by 5 × 10
repeated nested cross-validation with imputation, LASSO and tuning
inside each fold, and chosen by a pre-specified rule; the selected model
was frozen and evaluated once on a held-out set. A reference cohort with
the benchmark's criteria tested the original pipeline. As exploratory
extensions, first-day trajectory features, prior-diagnosis embeddings
and stacking were compared on the development set.

\textbf{Results.} In the reference cohort, the benchmark's RF reproduced
its training AUC of 0.999 but reached only 0.724--0.767 under nested
cross-validation, against 0.791 for the best model. In the landmark
cohort (2,557 patients, 26.6\% mortality), six configurations were
equivalent and L2-penalised LR was selected (nested-CV AUROC 0.801). Its
nested-CV AUPRC was 0.594 (prevalence 0.266). On the 512-patient
held-out set its AUROC was 0.750 (95\% CI 0.702--0.797), its AUPRC 0.533
(0.459--0.611) and its calibration slope 0.82. Red cell distribution
width, Charlson comorbidity index and age were the leading predictors.
Stacking raised AUROC by 0.005 over LR (95\% CI $-$0.001 to 0.011), and
trajectory features and diagnosis embeddings gave no gain.

\textbf{Conclusions.} Under leakage-free evaluation, LR matches tree
ensembles and the benchmark's RF advantage is not supported. The
landmark model can be applied at a defined time but needs local
recalibration before use.
